% Part: computability
% Chapter: machines-computations
% Section: halting-problem

\documentclass[../../../include/open-logic-section]{subfiles}

\begin{document}

\olfileid[hi]{tur}{und}{hal}
\olsection{विराम समस्या}

\begin{explain}
मान लें कि हमने ट्यूरिंग मशीन के वर्णनों का कोई परिगणन निश्चित किया है।
इस प्रकार प्रत्येक ट्यूरिंग मशीन को एक \emph{सूचकांक} मिलता है: ट्यूरिंग
मशीन के वर्णनों के परिगणन $M_1$, $M_2$, $M_3$, \dots{} में उसका स्थान।

हम जानते हैं कि ट्यूरिंग-असंगणनीय फलन अवश्य हैं: ट्यूरिंग मशीन के
वर्णनों का समुच्चय---और इसलिए ट्यूरिंग मशीनों का समुच्चय---!!{enumerable}
है, किंतु $\Nat$ से $\Nat$ तक के सभी फलनों का समुच्चय ऐसा नहीं है। फिर
भी हम असंगणनीय फलनों के विशिष्ट उदाहरण भी पा सकते हैं। ऐसा एक फलन विराम
फलन है।
\end{explain}

\begin{defn}[विराम फलन]
 \emph{विराम फलन}~$h$ को इस प्रकार परिभाषित करते हैं:
\[
h(e,n) =
\begin{cases}
  \text{0} & \text{यदि मशीन~$M_e$ आगत $n$ पर नहीं रुकती} \\
  \text{1} & \text{यदि मशीन~$M_e$ आगत $n$ पर रुकती है}
\end{cases}
\]
\end{defn}

\begin{defn}[विराम समस्या]
\emph{विराम समस्या}, किसी भी $e$, $n$ के लिए यह निर्धारित करने की समस्या
है कि ट्यूरिंग मशीन~$M_e$,~$n$ रेखाओं वाले आगत पर रुकती है या नहीं।
\end{defn}

\begin{explain}
हम सिद्ध करते हैं कि~$h$ ट्यूरिंग-संगणनीय नहीं है; इसके लिए हम एक
संबद्ध फलन~$s$ को ट्यूरिंग-असंगणनीय दिखाते हैं। यह प्रमाण इस तथ्य पर
निर्भर है कि ट्यूरिंग
मशीन द्वारा संगणित की जा सकने वाली प्रत्येक वस्तु की संगणना अनुशासित
ट्यूरिंग मशीन (\olref[mac][dis]{sec}) द्वारा की जा सकती है, और इस तथ्य
पर कि दो ट्यूरिंग मशीनों को जोड़कर एक अकेली मशीन बनाई जा सकती है
(\olref[mac][cmb]{sec})।
\end{explain}

\begin{defn} फलन~$s$ को इस प्रकार परिभाषित करते हैं:
\[
s(e) =
\begin{cases}
  \text{0} & \text{यदि मशीन~$M_e$ आगत $e$ पर नहीं रुकती} \\
  \text{1} & \text{यदि मशीन~$M_e$ आगत $e$ पर रुकती है}
\end{cases}
\]
\end{defn}

\begin{lem}
फलन~$s$ ट्यूरिंग-संगणनीय नहीं है।
\end{lem}

\begin{proof}
विरोधाभास के लिए मान लें कि फलन~$s$ ट्यूरिंग-संगणनीय है। तब कोई
ट्यूरिंग मशीन~$S$ होगी जो~$s$ की संगणना करती है। व्यापकता की हानि के बिना
मान सकते हैं कि $S$ जब रुकती है तब पहले खाने को पढ़ रही होती है
(अर्थात् वह अनुशासित है)। इस मशीन को दूसरी मशीन~$J$ से ``जोड़ा'' जा
सकता है, जो आगत~$0$ पर आरंभ किए जाने पर रुकती है (अर्थात् आरंभिक अवस्था
में पट्टी-अंत प्रतीक की दाईं ओर का खाना पढ़ते हुए यदि उसे $\TMblank$
मिलता है), और अन्यथा कभी न रुकते हुए दाईं ओर भटकती जाती है।
$S \concat J$, अर्थात् $S$ को~$J$ से जोड़कर बनी मशीन, एक ट्यूरिंग मशीन
है; अतः वह $M_e$ है, किसी~$e$ के लिए (अर्थात् परिगणन में कहीं आती है)।
$M_e$ को~$e$ संख्या
के $\TMstroke$ वाले आगत पर आरंभ करें। दो संभावनाएँ हैं: या तो $M_e$
रुकती है या नहीं रुकती।
\begin{enumerate}
\item मान लें कि $M_e$, $e$ संख्या के $\TMstroke$ वाले आगत पर रुकती
  है। तब $s(e)
  = 1$। अतः~$S$,~$e$ पर आरंभ होने पर पट्टी पर निर्गत के
  रूप में एक अकेले $\TMstroke$ के साथ रुकती है। फिर~$J$ पट्टी पर एक
  $\TMstroke$ के साथ आरंभ होती है। इस स्थिति में $J$ नहीं रुकती। किंतु
  $M_e$ ही मशीन $S
  \concat J$ है, इसलिए उसे ठीक वही करना चाहिए जो पहले
  $S$ और फिर $J$ करतीं (अर्थात् इस स्थिति में दाईं ओर भटकना और कभी न
  रुकना)। अतः $M_e$, $e$ संख्या के $\TMstroke$ वाले आगत पर नहीं रुक सकती।

\item अब मान लें कि $M_e$, $e$ संख्या के $\TMstroke$ वाले आगत पर नहीं
  रुकती। तब $s(e) = 0$, और~$S$, आगत~$e$ पर आरंभ होने पर रिक्त पट्टी के
  साथ रुकती है। रिक्त पट्टी पर आरंभ होने वाली~$J$ तुरंत रुकती है। फिर,
  $M_e$ वही करती है जो पहले $S$ और फिर~$J$ करतीं, अतः $M_e$ को $e$
  संख्या के $\TMstroke$ वाले आगत पर रुकना ही चाहिए।
\end{enumerate}
दोनों स्थितियों में हमें अपनी मान्यता से विरोधाभास मिलता है। इससे
दिखता है कि ऐसी ट्यूरिंग मशीन~$S$ नहीं हो सकती: $s$ ट्यूरिंग-संगणनीय
नहीं है।
\end{proof}

\begin{thm}[विराम समस्या की असमाधेयता]
\ollabel{thm:halting-problem} विराम समस्या असमाधेय है; अर्थात् फलन~$h$
ट्यूरिंग-संगणनीय नहीं है।
\end{thm}

\begin{proof}
मान लें कि $h$ ट्यूरिंग-संगणनीय है, मान लें किसी ट्यूरिंग मशीन~$H$ द्वारा।
हम~$H$ का उपयोग करके~$s$ की संगणना करने वाली ट्यूरिंग मशीन बना सकते हैं:
पहले आगत की एक प्रतिलिपि बनाएँ (जिसे एक~$\TMblank$ प्रतीक से पृथक किया
गया हो)। फिर आरंभ पर लौटकर~$H$ चलाएँ। स्पष्टतः हम पहली क्रिया करने वाली
मशीन बना सकते हैं (\cref{tur:mac:dis:prob:copier} देखें), और यदि $H$
का अस्तित्व होता तो उसे ऐसी प्रतिलिपिकार मशीन से ``जोड़कर'' नई मशीन पा
सकते थे जो निर्धारित करती कि $M_e$ आगत~$e$ पर रुकती है या नहीं, अर्थात्
~$s$ की संगणना करती। किंतु हम पहले ही दिखा चुके हैं कि ऐसी मशीन का
अस्तित्व नहीं हो सकता। अतः $h$ भी ट्यूरिंग-संगणनीय नहीं है।
\end{proof}

\begin{prob}
त्रि-विराम (3-विराम) समस्या ऐसी निर्णय-प्रक्रिया देने की समस्या है जो
निर्धारित करे कि मनमाने ढंग से चुनी गई ट्यूरिंग मशीन, अन्यथा रिक्त पट्टी
पर तीन $\TMstroke$ वाले आगत के लिए रुकती है या नहीं। सिद्ध करें कि
3-विराम समस्या असमाधेय है।
\end{prob}

\begin{prob}
दिखाएँ कि यदि विराम समस्या उन ट्यूरिंग मशीन और आगत युग्मों $M_e$ और~$n$
के लिए समाधेय है जहाँ $e \neq n$, तो वह उन स्थितियों के लिए भी समाधेय
है जहाँ $e = n$।
\end{prob}

\begin{prob}
हमने सिद्ध किया कि यदि आगत कोई संख्या~$e$ है, जो सभी ट्यूरिंग मशीनों के
परिगणन द्वारा ट्यूरिंग मशीन~$M_e$ को पहचानती है, तो विराम समस्या असमाधेय
है। यदि हम \olref[tur][und][enu]{sec} के ट्यूरिंग मशीन-वर्णन को सीधे
आगत के रूप में लेने दें तो क्या होगा? क्या ऐसी ट्यूरिंग मशीन हो सकती है
जो विराम समस्या का निर्णय करती हो किंतु आगत के रूप में सूचकांकों के
स्थान पर ट्यूरिंग मशीनों के वर्णन लेती हो? समझाएँ कि क्यों या क्यों नहीं।
\end{prob}

\begin{prob} दिखाएँ कि निम्न प्रकार परिभाषित \emph{आंशिक} फलन~$s'$
  \[
  s'(e) =
  \begin{cases}
    \text{1} & \text{यदि मशीन~$M_e$ आगत $e$ पर रुकती है}\\
    \text{अपरिभाषित} & \text{यदि मशीन~$M_e$ आगत $e$ पर नहीं रुकती}
  \end{cases}
  \]
  ट्यूरिंग-संगणनीय \emph{है}।
\end{prob}

\end{document}
